%======================================================================
% INTERFAZ PRINCIPAL
%======================================================================

\chapter{Interfaz Principal}

La consola está organizada en módulos funcionales que reflejan las operaciones empresariales del sistema. Este capítulo describe la estructura de navegación y los componentes principales de la interfaz.

\section{Arquitectura Angular Enterprise}

La arquitectura de componentes standalone representa un cambio paradigmático en Angular:

Los Componentes Standalone son una característica de Angular que permite crear componentes sin necesidad de NgModules, simplificando la arquitectura y mejorando el rendimiento de las aplicaciones.

Los componentes standalone eliminan la necesidad de módulos intermedios, reduciendo el tiempo de carga y mejorando el tree-shaking en producción.

Las ventajas implementadas incluyen:

\begin{itemize}
\item \textbf{Reactividad Instantánea}: Uso de Signals para cambio detection eficiente sin zone.js overhead.
\item \textbf{Tree-shaking}: Eliminación de código no utilizado en producción.
\item \textbf{Lazy Loading Natural}: Carga de componentes bajo demanda sin configuración adicional.
\item \textbf{Testabilidad}: Componentes más fáciles de unit testear de forma isolado.
\end{itemize}

La estructura de navegación modular permite escalar la aplicación sin comprometer el rendimiento, cargando solo los módulos necesarios.

\section{Módulos de Navegación Principal}

La consola incluye los siguientes módulos de navegación:

\begin{itemize}
\item \textbf{Login}: Punto de acceso seguro con autenticación robusta
\item \textbf{Principal}: Dashboard central con métricas en tiempo real
\item \textbf{Redes Informáticas}: Gestión y monitoreo de infraestructura de red
\item \textbf{Herramientas}: Utilidades administrativas y de diagnóstico
\item \textbf{Aplicaciones}: Gestión de aplicaciones del ecosistema
\item \textbf{Seguridad}: Centro de control de seguridad y auditoría
\item \textbf{Opciones}: Configuración del sistema y preferencias
\end{itemize}

\section{Dashboard Principal}

El dashboard principal proporciona una vista consolidada del estado del sistema:

\begin{itemize}
\item \textbf{Métricas en Tiempo Real}: Uso de CPU, memoria, conexiones activas.
\item \textbf{Estado de Servicios}: Disponibilidad de cada componente del ecosistema.
\item \textbf{Notificaciones}: Alertas y eventos importantes.
\item \textbf{Accesos Rápidos}: Acciones frecuentes para administradores.
\end{itemize}

\clearpage
